\section{One source tree for every normal cut}
\label{sec:templates}

The cap graph changes when a source edge weight becomes positive. It would
be wasteful to regard each zero/nonzero support as an unrelated scheduling
problem. The source one-skeleton already supplies one schedule template for
all of them.

\begin{definition}[Boundary-rooted source forest]
Choose a forest $T_0$ in the source vertex graph such that every interior
vertex has a unique parent edge and every component has one source boundary
vertex as its root. There is no parent edge for a boundary vertex. Record the
interior vertices in a parent-before-child order. Such a forest is produced
by a multi-source breadth-first search starting at all boundary vertices.
\end{definition}

The term forest refers to the actual source vertices. Contracting all
boundary roots to one formal vertex makes its interior portion a rooted
tree for the scheduling argument. This is still only graph bookkeeping.

\begin{theorem}[Universal source-tree specialization]
\label{thm:template}
Fix $T_0$ independently of normal coordinates. For every admissible $x$,
process its interior vertices in the recorded order. For a vertex $v$ with
parent edge $e$, choose
\begin{enumerate}
\item the unique exceptional gap of $e$, when $w_e=0$;
\item the extreme gap of $e$ adjacent to $v$, when $w_e>0$.
\end{enumerate}
Cancel that two-handle with the three-handle of $v$. This is a complete
valid trace for the relative operation of \cref{thm:top}. The same source
template works for all admissible coordinate supports and all scales.
\end{theorem}

\begin{proof}
An edge incident to an interior vertex is not a boundary edge, so the
required exceptional gap exists. If its weight is zero, its other endpoint
is the parent, which is a boundary vertex or has already been removed. It
therefore has exactly one live three-handle contact, namely $v$. If its
weight is positive, the child-adjacent extreme gap has its other endpoint on
the new cut boundary; it is immediately eligible, regardless of the
parent's status.

Distinct vertices use distinct parent edges of the source forest. Two
different choices can consequently never request the same exceptional
two-handle, including after a support change replaces one zero gap by two
extreme gaps. A parent edge cannot be an interior loop. Every interior
vertex is processed once, and \cref{lem:relative-move} applies at every
step. The construction makes no assumption on coordinate size and includes
$x=0$.
\end{proof}

\begin{corollary}[Support and scale invariance]
\label{cor:support}
The cap graph, considered with oriented extreme-gap labels, depends on $x$
only through the predicates $w_e=0$. Positive scaling $x\mapsto kx$ leaves
the source-tree choices and the corresponding dense two-handle IDs
unchanged. A specialized trace has the same number of moves and source-ID
bits at every positive scale.
\end{corollary}

\begin{proof}
The cap formula only distinguishes zero from positive weight. In the latter
case the numerical label $w_e$ changes but the order of the two extreme
gaps does not. Source vertices, source edge order, and the parent forest
are fixed. The trace records source handle IDs, while its input fingerprint
has fixed encoded length.
\end{proof}

The corollary concerns the added top trace. The complete source certificate
still contains exact coordinates and affine endpoints and can grow with
$b_x$. It would be incorrect to infer constant total proof size from constant
top-trace size.

\subsection{A linear queue and an independent trace checker}

The default implementation uses a queue rather than requiring the reusable
template. Initialize every edge with one live endpoint. When a vertex is
deleted, revisit its incident edges and enqueue newly eligible ones.
Duplicate queue entries are harmless: a popped edge is checked against the
current live sets before use. Each interior loop is stored once in the
vertex adjacency list but retains both endpoint occurrences in the edge
record.

\begin{proposition}[Local cost and trace size]
\label{prop:cost}
Given an authenticated dense-ID cap packet with $m$ edges and $q$ interior
vertices, a complete rooted schedule and the surviving edge list can be
produced in $O(m+q)$ word operations and $O(m+q)$ space. The supplied trace
can be checked in the same structural bound. For the source cores in
\eqref{eq:core-bounds}, the additional proof takes
$O(t\log(t+2))$ bits. Source validation, coordinate reading, rank address
reconstruction, and fingerprinting remain polynomial in
$t\log(t+2)+b_x$.
\end{proposition}

\begin{proof}
The initialization and final dense-ID survivor scan visit each edge once.
There are at most $2m$ vertex--edge incidences. Each is revisited only when
its vertex is deleted. Each visit performs constant work on an edge's two
endpoint occurrences and can enqueue it at most once. The initial queue
also has at most $m$ entries. Thus all queue pops and endpoint checks have
linear total count. A checker reads each of the $q$ pairs, rejects reuse,
and tests its at most two live contacts. It finally compares the exact
survivor list.

All original handle IDs are $O(t)$, so the $q$ pairs and at most $m$
survivors use $O(t\log(t+2))$ bits. The template also consists of $O(t)$
vertex, parent, and source-edge labels of this length. Reconstructing the
source data requires a polynomial number of combinatorial operations and
arithmetic on the input integers. No loop ranges over the numerical normal
weight, and the fingerprint uses hexadecimal coordinate encodings rather
than decimal conversions with implementation-specific digit limits.
\end{proof}

The public template functions validate their source arguments rather than
assuming a previously trusted cache. Thus their current complete-call cost
includes source validation even when the combinatorial template is reused.
A future authenticated session object could remove duplicate validation
work while preserving the source binding. Such caching is a separate
implementation improvement, not part of the linear queue claim.

\subsection{How the source orientation is authenticated}

A triangulation edge can appear in many local tetrahedra, and its two global
endpoint vertices can coincide. The canonical orientation is determined by
the signed source edge union, not by sorting vertex labels. The producer
uses the union-find root's own oriented local representative. The verifier
uses a different reconstruction: it scans \emph{every} local occurrence,
reverses its endpoint pair according to the recorded orientation parity,
and demands agreement.

This matters in practice. The general source validator keeps a convenient
last-seen endpoint pair that is sufficient for link counting but need not
be in the canonical orientation. Using that pair to decide which extreme
gap has rank zero can reverse the cap formula. The new proof checker does
not trust that shortcut. It reconstructs the complete expected exceptional
rank inventory before interpreting any move.

\subsection{The exact scope of uniformity}

The template is universal for one fixed source triangulation. A subsequent
hierarchy state may have a different triangulation or a differently
represented marked handle structure. The theorem does not identify those
states, transport their patterns, or establish a bounded number of source
rebuilds. Its useful implication is more local: within a fixed source,
top-handle elimination introduces no branch over normal supports and no
normal-weight-dependent number of moves.
